\section{Statement, proof contract, and logical route}
\label{cp:sec-statement}

\subsection{Spherical scissors congruence}

Let \(S^3\subset\R^4\) be the unit sphere.  A spherical polytope is a
finite union of geodesic simplices contained in open hemispheres.  The
spherical scissors group \(\Pfull(S^3)\) is generated by symbols \([P]\),
one for each such polytope, subject to
\[
 [P]=[P_1]+\cdots+[P_m]
\]
when the \(P_i\) have disjoint interiors and union \(P\), and
\([P]=[gP]\) for every spherical isometry \(g\), including reflections.  This
is the unoriented geometric scissors group used in the conclusion; the
orientation of a complex quadric introduced later is auxiliary data, carried
through the parity line in Input~\textup{(E6)}.  Equality in this group is
the stable algebraic form of equidecomposability.  The cancellation theorem
listed in Input~\textup{(E8)} below turns the equality obtained at the end
into a finite equidecomposition.

Label the vertices of a nondegenerate spherical tetrahedron by
\(0,1,2,3\), and select the opposite edge pair \(01,23\).  The elementary
Regge operation fixes those two edge lengths and replaces the remaining
four lengths \(\ell_1,\ldots,\ell_4\) by
\begin{equation}
 \ell_i^R=\sigma_\ell-\ell_i,
 \qquad
 \sigma_\ell=\frac12\sum_{j=1}^4\ell_j.
 \label{cp:eq-regge-lengths}
\end{equation}
Input~\textup{(E1)} says that the same transformation holds for the four
moving dihedral angles and that the two tetrahedra have equal volume.

\begin{theorem}[Conditional spherical Regge scissors theorem]
\label{cp:thm-main}
Assume Inputs \textup{(E1)--(E8)} below, including the relative motivic
bridge in Inputs \textup{(E2)--(E5)}, for the two universal rational charts
of Section~\ref{cp:sec-phase-charts}.  If \(T\subset S^3\) is a
nondegenerate spherical tetrahedron and \(T^R\) is a nondegenerate
elementary Regge mate, then
\[
 [T]=[T^R]\qquad\text{in }\Pfull(S^3).
\]
Consequently \(T\) and \(T^R\) admit a finite spherical
equidecomposition.
\end{theorem}

No algebraicity hypothesis is imposed on the angles or vertices.  The
qualifier \emph{conditional} is essential: this manuscript proves the
implication from the following precisely typed inputs; it does not claim
that the relative motivic bridge has already been proved in the cited
literature.

\subsection{The external theorem package}
\label{cp:subsec-external-contract}

We separate published theorems from the additional comparison statements
needed to use them in a family.  This separation is also the interface used
by the Lean formalization.

\begin{enumerate}[label=\textup{(E\arabic*)},leftmargin=3.1em]
\item \textbf{Regge geometry.}
\label{cp:input-E1}
For a genuine spherical Regge pair, the same linear involution acts on the
four moving edge lengths and on the four moving dihedral angles, while the
chosen opposite pair is fixed; the two tetrahedra have equal volume.  This
is imported from Akopyan--Izmestiev \cite{AkopyanIzmestiev2019}.

\item \textbf{Relative quadric coefficients.}
\label{cp:input-E2}
For each rational chart \(U\) constructed in Section
\ref{cp:sec-phase-charts}, the oriented quadric families attached to
\(G\) and \(G^R\) determine elements
\[
 c_U(G),c_U(G^R)\in A_{02}(U)
\]
of a rational vector space, and there is a linear reduced coproduct
\[
 \overline\Delta_U:A_{02}(U)\longrightarrow A_{11}(U).
\]
These data are functorial under pullback to a complex point.  Their intended
construction is the framed relative face motive described in Section
\ref{cp:sec-motivic-dictionary}; existence and functoriality over this
positive-dimensional base are input, not an internal consequence of the
matrix calculation.

\item \textbf{Three-weight and primitive comparison.}
\label{cp:input-E3}
There is a rational vector space \(H_U\), an identification
\[
 \ker\overline\Delta_U\xrightarrow{\sim}H_U,
 \qquad H_U=H^1(U,\Q(2)),
\]
and \(H_U=0\).  The intended proof combines a three-step
Artin--Tate filtration, a heart-to-derived \(\Ext^1\) comparison,
weight-two birational invariance, and Borel's rank theorem.  Garkusha
\cite{Garkusha2023} and Borel \cite{Borel1977} supply the cited
birational and number-field ingredients.  The assertion that the particular
relative face objects lie in a strict heart with this primitive kernel is
part of the input.

\item \textbf{Relative coproduct formula and descent.}
\label{cp:input-E4}
The coproduct of a framed quadric coefficient is the sum of its six
oriented length--angle channels,
\[
 \overline\Delta_U c_U(G)=\sum_e L_e\otimes A_e,
\]
and similarly for \(G^R\).  Five channels are ordinary Tate channels and
the fixed \(p\)-channel may be a sign-Artin channel.  Pullback to a splitting
cover, the normalized twisted trace, and the displayed formula are
compatible, so equality after pullback descends to \(U\).  Pullback to every
complex point commutes with the reduced coproduct and identifies the pulled
back six channels with Goncharov's field-level coproduct; in particular, a
universal zero-coproduct defect specializes to a primitive quadric-scissors
class.  This input also
includes the algebraic recovery of the length phases and the Zariski-density
step which extends the real Regge relations to that splitting cover.
Goncharov's
field-level coproduct calculation \cite{Goncharov1999} motivates this
input; its exact relative and Artin-trace form is not proved here.

\item \textbf{Fibre comparison.}
\label{cp:input-E5}
For every complex point \(s:\Spec\C\to U\), specialization of relative
coefficients and passage to the endpoint are linear maps.  Let
\(\delta_s^\vee\) be the dual quadric-scissors defect.  If
\(\delta_U^\vee=c_U(G)-c_U(G^R)\), then the endpoint obtained from
\(s^*\delta_U^\vee\) is exactly
\[
 c_G(\delta_s^\vee)
 \in \gr_2^\gamma K_3(\C)\otimes\Q,
\]
with the same signs, projective duality, frames, and Tate twist as in
Goncharov's definition.  This is an equality of classes, not merely of
regulators.  The family-to-fibre identification is an explicit input.

\item \textbf{Goncharov's field-level injections.}
\label{cp:input-E6}
The canonical map from primitive weight-two quadric-scissors classes to
\(\gr_2^\gamma K_3(\C)\otimes\Q\) is injective, and the spherical map from
the reduced spherical scissors group to the spherical part of the quadric
group is injective.  We use the projective-duality involution as well, with
the explicit generator compatibility taking the spherical defect represented
by \(Q_{G^{-1}}\) to the dual defect represented by \(Q_G\).  These
are the precise consequences of Goncharov's Theorems~4.13 and~4.14 used
below \cite{Goncharov1999}.

\item \textbf{Spherical presentation comparison.}
\label{cp:input-E7}
The reduced geometric spherical scissors group is the spherical group on
which Input~\textup{(E6)} is defined.  After rationalization the comparison is
an isomorphism and carries the particular reduced Regge defect to the
particular Goncharov spherical defect used by Input~\textup{(E6)}.  This is
the comparison with the
coinvariant configuration presentation supplied by Dupont's
Theorem~7.14 \cite{Dupont2001}.

\item \textbf{Elementary spherical scissors facts.}
\label{cp:input-E8}
The groups \(\Pfull(S^2)\) and \(\Pfull(S^3)\) are uniquely divisible.  The
reduction map to \(\Pred(S^3)=\operatorname{coker}\Sigma\) is additive and
exact at \(\Pfull(S^3)\): a class has zero reduction exactly when it is a
suspension.  The canonical rationalization map on \(\Pred(S^3)\) is
injective.  Area and volume are additive, area identifies
\(\Pfull(S^2)\) with \(\R\), and equality in
\(\Pfull(S^3)\) implies finite equidecomposability by Zylev cancellation
(with Gerling's orientation refinement).  We also take as input the
normalization
\[
 \Vol(\Sigma Q)=\frac\pi2\Area(Q)
\]
for spherical suspension.  The scissors statements are imported in the
form recorded by Dupont \cite{Dupont2001}; the normalization is the
warped-product volume calculation displayed in Section
\ref{cp:sec-specialization}.
\end{enumerate}

Inputs \textup{(E2)--(E5)} are deliberately stronger and more explicit
than the slogan \emph{motivic formalism is standard}.  The cited sources
develop the field-level constructions and the ingredients suggested above,
but this manuscript does not supply a reference proving their exact
simultaneous relative form.  Until that bridge is constructed, the result
of this paper is the conditional theorem stated above.

\subsection{What Lean checks}

The accompanying Lean development does not define Voevodsky motives,
motivic cohomology, or algebraic \(K\)-theory.  Instead it represents
Inputs \textup{(E1)--(E8)} by explicit structures and proves, without
additional axioms, the implication from an inhabitant of those structures
to the final scissors equality.  It also checks the finite algebra that is
specific to Regge symmetry: the involution and orthogonality of the Regge
matrix, its monomial phase action, the coordinate identities and covering
criterion for the two orientation charts, the fixed quadratic deck
transformation, the tensor cancellation, the
primitive-kernel deduction, and the final diagram chase.  The exact
manuscript-to-declaration map is in \texttt{formal/MANUSCRIPT\_MAP.md}.
The scheme-level principal-open identifications and finite-etaleness
statements are standard external algebraic-geometry packaging and are marked
as mixed rather than kernel-checked items in that map.

This distinction is important.  A Lean theorem universally quantified over
an external package is a verified conditional theorem; it is not a formal
construction of that package and does not certify the cited analytic,
algebro-geometric, or motivic results.

\subsection{Route of the conditional proof}

The proof has seven steps.

\begin{enumerate}[label=\arabic*.,leftmargin=2.1em]
\item Replace each dihedral angle \(\theta\) by its phase
\(u=e^{\mathrm i\theta}\).  The additive Regge formula becomes the monomial
substitution \(u_i^R=\tau/u_i\).

\item Adjoin the orientation root \(w=\sqrt{\det G}\).  The substitutions
\(r=w-sx\) and \(r_+=w+sx\) rationalize the two necessary charts.  Every
physical pair gives a complex point of one of the resulting opens \(U\).

\item Apply Input~\textup{(E2)} to form
\(\delta_U^\vee=c_U(G)-c_U(G^R)\).  Section
\ref{cp:sec-motivic-dictionary} explains the relative face motive which this
symbol is intended to denote.

\item Use Input~\textup{(E4)} to translate its coproduct into six
length--angle tensors.  The two fixed tensors cancel directly; orthogonality
of the Regge matrix cancels the other four.  The descent clause handles the
single possible sign-Artin channel.

\item Thus \(\delta_U^\vee\in\ker\overline\Delta_U\).  Input
\textup{(E3)} identifies this kernel with the zero group
\(H^1(U,\Q(2))\), so \(\delta_U^\vee=0\).

\item Pull back to the point of the original tetrahedron.  Inputs
\textup{(E5)--(E8)} and the stated comparison and rationalization injections
give zero in the reduced spherical scissors group.

\item The unreduced difference is a suspension.  Regge volume equality,
the suspension-volume formula, and Input~\textup{(E8)} kill that suspension
and produce a finite equidecomposition.
\end{enumerate}

Section~\ref{cp:sec-phase-charts} starts with the fully explicit algebra in
Steps~1--2.
